\subsection{معماری و پیاده‌سازی}
این ماژول در فایل \lr{Codes/framework/compression\_engine/armor\_pruner.py} تحت دو کلاس اصلی طراحی و در رجیستری فشرده‌سازی با نام‌های \lr{armor\_pruner}، \lr{armor} و \lr{armor\_compressor} ثبت شده است:
\begin{itemize}
    \item \textbf{کلاس \lr{ARMORLinear}:} جایگزین لایه‌های خطی استاندارد (\lr{nn.Linear} و \lr{Conv1D}) می‌شود. در فرآیند گذر پیشرو (\lr{Forward Pass})، محاسبات به صورت دو شاخه موازی انجام می‌پذیرد:
    \begin{equation}
    \mathbf{y} = \mathbf{x} (\mathbf{W}_{\text{sparse}} \odot \mathbf{M}_{2:4})^\top + \frac{\alpha}{r} (\mathbf{x} \mathbf{A}^\top) \mathbf{B}^\top + \mathbf{b}
    \end{equation}
    هسته تنک ماتریس پایه ثابت شده و شاخه کم‌مرتبه قابلیت تنظیم دقیق گرادیانی را در صورت نیاز حفظ می‌کند. ماسک تنک به عنوان یک بافر غیردائمی (\lr{non-persistent buffer}) ذخیره شده تا در اندازه فایل مدل سربار ایجاد نکند.
    \item \textbf{کلاس \lr{ARMORPruner}:} کنترل‌کننده فرآیند هرس که با اعمال فیلترهای معماری (\lr{should\_compress\_linear}) لایه‌های هدف را شناسایی کرده، داده‌های کالیبراسیون را استخراج می‌نماید، برآورد حساسیت را انجام داده و با ارزیابی آستانه خطای فروبنیوس (\lr{max\_error})، در صورت فراتر رفتن خطا، مسیر بازگشت ایمن به وزن‌های اولیه (\lr{fallback}) را اجرا می‌کند.
\end{itemize}

\subsection{محدودیت‌ها و تحلیل شکست}
تحلیل رفتار و سازوکار فشرده‌سازی این ماژول محدودیت‌های زیر را برجسته می‌سازد:
\begin{itemize}
    \item \textbf{سربار حافظه در مرحله مقداردهی کم‌مرتبه:} اگرچه ماتریس تنک ساختاریافته حجم محاسبات لایه را نصف می‌کند، اما افزودن شاخه جبرانی \lr{LoRA} برای رتبه‌های بالا ($r \ge 32$) ممکن است در مدل‌های دیکودری بزرگ به افزایش تعداد کل پارامترها منتهی شود (مانند رخداد ثبت‌شده در بنچ‌مارک \lr{GPT-2} تحت نام \lr{model grew}).
    \item \textbf{نیاز به پشتیبانی سخت‌افزاری \lr{Sparse Core}:} شتاب‌بخشی خطی ساختار \lr{2:4} تنها در پردازنده‌های گرافیکی دارای معماری \lr{NVIDIA Ampere} و جدیدتر (از طریق دستورات \lr{mma.sp}) محقق می‌شود؛ در پردازنده‌های قدیمی‌تر، هزینه نگهداری ماسک و شاخه‌های موازی می‌تواند باعث افت جزئی توان پردازش گردد.
\end{itemize}
